\subsection{Augmentation kernels for a linear category}

\begin{proposition}\label{C:linear-augmentation-Tor}
Let $C$ be a finite category on objects $0,1,2$, with no
backward arrows and with trivial endomorphism monoids at
$0$ and $2$.  Suppose the three forward hom-sets are nonempty.
Put
\[
 N=\operatorname{End}_C(1),\quad R=\mathbb QN,\quad
 V=\mathbb QC(0,1),\quad W=\mathbb QC(1,2),\quad
 U=\mathbb QC(0,2).
\]
Here $V$ is a left $R$-module, $W$ is a right $R$-module,
and all three permutation spaces have their usual
augmentation to $\mathbb Q$.  Write $\bar V,\bar W,\bar U$
for the respective augmentation kernels.  Then, naturally,
\begin{equation}\label{C:linear-augmentation-high}
 H_n(BC;\mathbb Q)\cong
 \operatorname{Tor}^{R}_{n-2}(\bar W,\bar V)
 \qquad(n\geq3).
\end{equation}
There is also an exact sequence
\[
 0\longrightarrow H_2(BC;\mathbb Q)
 \longrightarrow\bar W\otimes_R\bar V
 \xrightarrow{\bar\mu}\bar U
 \longrightarrow H_1(BC;\mathbb Q)
 \longrightarrow0,
\]
where $\bar\mu$ is induced by the actual composition
$C(1,2)\times C(0,1)\to C(0,2)$.
\end{proposition}
\begin{proof}
Let $G$ be the full subcategory on $0,2$.  Its nerve is
the graph with those two vertices and one edge for each
element of $C(0,2)$.  It is connected.  Let $L_*$ be its
rational normalized chain complex, regarded as a
subcomplex of the normalized chains $C_*(NC;\mathbb Q)$.
It has no terms above degree one,
$H_0(L_*)=\mathbb Q$, and $H_1(L_*)=\bar U$.

Describe the quotient complex by the number of strict
changes of object.  With zero changes, only the middle
corner bar $B_*(\mathbb Q,R,\mathbb Q)$ remains.
With one change, its two summands are
$B_*(\mathbb Q,R,V)$ and $B_*(W,R,\mathbb Q)$.
With two changes, it is $B_*(W,R,V)$.  The one-change
and two-change terms have total-degree shifts one and
two.  The horizontal faces are the two augmentation
maps, with the ordinary alternating signs.  The face
composing the two cross arrows has its value in $U$;
it is zero in the quotient.  This description follows
by retaining the ordered string of middle endomorphisms
between the cross arrows in each nondegenerate simplex.
It includes the empty string and all its boundary faces.

Consequently the quotient complex is the total two-sided
bar for the two complexes
\[
 [\,W\xrightarrow{\varepsilon_W}\mathbb Q\,]
 \quad\hbox{and}\quad
 [\,V\xrightarrow{\varepsilon_V}\mathbb Q\,],
\]
where $W,V$ occupy degree one and $\mathbb Q$ degree zero.
Their augmentations are onto, so these complexes are
quasi-isomorphic to $\bar W[1]$ and $\bar V[1]$.
The two-sided bar computes their derived tensor product
over $R$.  These quasi-isomorphisms can also be used
without an abstract derived comparison: in each fixed
internal bar degree all other tensor factors are vector
spaces over $\mathbb Q$, hence preserve exact sequences;
the first-quadrant total-complex comparison then preserves
quasi-isomorphisms.  Thus the quotient has homology
\[
 H_n\bigl(C_*(NC;\mathbb Q)/L_*\bigr)
 \cong\operatorname{Tor}_{n-2}^R(\bar W,\bar V),
\]
with the right side understood as zero for $n<2$.

The long exact sequence for $L_*$ gives
\eqref{C:linear-augmentation-high} in degrees at least
three and gives the displayed low-degree exact sequence.
Its connecting map is composition: on a tensor
$(w-w_0)\otimes(v-v_0)$, the two remaining augmentation
faces cancel, and the direct-edge face is
\[
 wv-wv_0-w_0v+w_0v_0.
\]
This lies in $\bar U$, respects the $R$-balancing
relations, and is precisely $\bar\mu$, up to the
irrelevant simultaneous sign from the chain convention.
\end{proof}

\begin{corollary}\label{C:semisimple-ends-middle-four}
Let $C$ have three linearly ordered objects and nonempty
forward hom-sets.  Suppose its two endpoint monoid
algebras are semisimple over $\mathbb Q$.  If its middle
monoid has at most four elements, or has a two-sided zero
and at most five elements, bounded rational homology of
$BC$ implies
\[
 H_4(BC;\mathbb Q)=0.
\]
In particular these categories cannot have type
$\mathbb CP^2$, at any cross-hom cardinalities.
\end{corollary}
\begin{proof}
Collapse each endpoint monoid to the trivial monoid,
replacing every incident hom-set by its actual orbit set.
The orbit-collapse functors and the exact-coinvariant
comparison of
Corollary~\ref{C:semisimple-ends-binary-homs}
give a rational homology equivalence.  The same
strict-object bar argument applies with arbitrary
nonempty hom-set sizes: semisimplicity makes endpoint
coinvariants exact on every associated-graded bar.
We may apply Proposition~\ref{C:linear-augmentation-Tor}
to this resulting category.

If $|N|\leq4$, adjoin a new absorbing zero $z$ to $N$,
even if $N$ already has another zero.  The contracted
algebra $\mathbb Q(N\sqcup\{z\})/\mathbb Qz$ is
exactly $\mathbb QN$, of dimension at most four.
Theorem~\ref{B:contracted-four-Tor} says that a nonzero
degree-two Tor group for the displayed augmentation
kernels forces nonzero Tor in arbitrarily high degrees.
Equation~\eqref{C:linear-augmentation-high} would then
contradict bounded rational homology of $BC$.

If $N$ has a two-sided zero and at most five elements,
its algebra is the product of the semisimple zero factor
$\mathbb Qz$ and the contracted algebra of dimension at
most four.  Positive Tor is contributed only by the
contracted factor, so the same theorem gives the same
contradiction.  Thus degree-two Tor is zero, and the
degree-four case of
\eqref{C:linear-augmentation-high} proves the conclusion.
\end{proof}

\begin{proposition}[Collapsing the long hom-set]
\label{C:linear-long-hom-collapse}
Under the hypotheses of
Proposition~\ref{C:linear-augmentation-Tor}, let $D$
retain every object and arrow of $C$ except that all
arrows $0\to2$ are identified to a single arrow.
Composition into that hom-set is the unique possible
map.  This defines a category and an actual functor
$C\to D$.  If $BC$ is simply connected and
$u=|C(0,2)|$, then
\[
 BD\simeq BC\vee\bigvee^{u-1}S^2.
\]
In particular, if $BC\simeq\mathbb CP^2$, the
triangle-compression monoid of order
$|N|+|C(0,1)|\,|C(1,2)|$ realizes
$\mathbb CP^2\vee\bigvee^{u-1}S^2$.
\end{proposition}
\begin{proof}
The only nondegenerate simplices containing a direct
$0\to2$ arrow are those edges themselves: the endpoint
corners are trivial.  Thus $ND$ is the simplicial pushout
of $NC$ and the single-arrow nerve along $NG$.
The inclusion $NG\hookrightarrow NC$ is a simplicial
subcomplex inclusion, so realization gives its homotopy
pushout.  The single-arrow nerve is contractible.
The graph $BG$ is a wedge of $u-1$ circles, and its
map into the simply connected space $BC$ is
null-homotopic.  The pushout is therefore
$BC\vee\Sigma BG$, which is the stated wedge.

For completeness the compression monoid is the disjoint
union $N\sqcup(S\times T)$, where
$S=C(0,1)$ and $T=C(1,2)$.  Its old multiplication is
that of $N$, and its other products are
\[
 (s,t)(s',t')=(s,t'),\qquad
 n(s,t)=(ns,t),\qquad (s,t)n=(s,tn).
\]
The identity is the identity of $N$.  The action laws
give associativity in every mixed case.  The triangle-compression comparison of
Theorem~\ref{C:triangle-band-compression} identifies its
classifying space with $BD$; here the long hom-set is
exactly a singleton, so all hypotheses of that comparison
hold.  Its arrow count is the displayed sum.
\end{proof}
